% !TeX root = ../papers/nc-p3.tex
\section{通常Bell曲面の単純点と反転なし衝突点に至る全窓}
\label{sec:split-surface-all-fibers}
\begin{theorem}\label{thm:split-surface-all-fibers}
通常Bell曲面の全ての単純スペクトルの閉軌道、および重複曲面の半単純・反転なしの
全閉軌道を取る。これらの全流入ファイバーで、原点を除く無印同型類は
$\mathbb P(1^4,2^3)$ の補正Pauli四元群による商である。
全点が標準Koszul左右AS正則で、階数スキームで滑らかであり、十五次元成分だけが通る。
各非零軌道の半安定境界は指定閉軌道のみである。異なる閉軌道上の類は重ならない。
追加単純表示の滑らかな60次元成分では、単純スペクトルの全商ファイバーが閉軌道一つだけになる。

通常表示では、非零類のうち周期が2を割る軌跡の既約成分は全てのファイバーで
同型類次元 $4,2,2,2$ である。定理\ref{thm:split-surface-periods}の固有空間式が
各点での必要十分条件で、それ以外の類は最小周期4である。
閉軌道が最小周期2ならこの四成分上の全点は最小周期2である。
残る反転なし衝突の $f(t)^2=1$ の閉軌道では、
非零類で最小周期1の軌跡は次元 $2,1,1,0,0$ の五つの既約成分を持つ。
\end{theorem}
\begin{proof}[パラメータ全域での階数]
通常四表示は基底パラメータの符号・逆数・交換の持上げで $h=1$ の表示へ移せる。
衝突の反転なしは $u=v=t$, $t^4\ne1$ の $h=1$ を使う。法線文字を0,1、他を2,3とする。
二つの切断の左右核から取る全接行列を、まず法線数 $n$ で分ける。
さらに定理\ref{thm:split-surface-periods}の符号 $Z$ と補正ビット反転 $X$ を使う。
前者の符号は語の総ビットの偶奇 $p$ で定まり、後者の固有値は
$n=2$ なら $\epsilon=\pm1$、$n=3$ なら $\epsilon\iota$ である。
$X$ は各語を全補語へ送り、法線位置で $(\iota,\iota,-\iota,-\iota)$ の積を掛ける。
二つの隣接切断の行・列核基底は、四つの異色対と二つの純関係ピボットを用いて選べる。
それらへの射影が全域で同型であることは可逆分解と純行列の単項形による。

\path{computations/verify_split_surface_all_fibers.py} はこの核の全係数の分母が
$D=uv(u^2-1)(v^2-1)$ の因子だけであることを調べる。
両核を掛けた各接方程式で、各語を一方のビットが0の代表と全補語の固有和に制限する。
$(p,\epsilon)=(0,+),(0,-),(1,+),(1,-)$ の順に、各ブロックの階数の下界は次の通りである。
\[
\begin{array}{c|c|c}
 &\text{各ブロックの列数}&\text{階数下界}\\\hline
 n=2 &24&18,16,16,16\\
 n=3 &16&11,11,11,11\\
 n=3\text{のせん断作用}&4&4,4,4,4
\end{array}
\]
下二行では $D$ だけを逆にする。第一行は
$D(u^2-v^2)(u^2v^2-1)$ と $D(u+v)(uv-1)(u^2+1)$ の各主開集合上で別々に検算する。
後者は上記の全衝突 $u=v=t$ を含む。実際の証明書は、各行の単元分母を払って単元ピボットで消去し、
残る必要階数の全小行列式と $1-zD'$ が単位イデアルを生成することを有理係数で確認するもの
（$D'$ は当該開集合の式）である。せん断行列には各位置の第一文字だけへのせん断を
補語と固有和にして使う。全16本のせん断空間の各指標をこれで尽くす。

$n=4$ の純行列の接核は三次元である。実際、通常表示の法線の純二次関係は
四つのPauli関係でパラメータに依存せず、定理\ref{thm:split-surface-full-fibers}の一般値をそのまま使える。
軸交換の自己同型は全パラメータでの有理式で与えられ、分母は $D$ の因子だけなので
$n=0,1$ にも $n=4,3$ の次元上界が成立する。従ってアフィン接次元は高々
$3+20+30+20+3=76$。全点を通る既約なアフィン76次元成分があるから全ての上界は等号であり、
局所環は正則である。滑らかな一成分点という不変な開条件は流入窓にも移る。
\end{proof}
\begin{proof}[横断表示と周期層]
単純点で分解の向きを保つ安定群は法線倍率と補正Pauli群である。
$u^2=-1$ などで交代成分の安定群が増えても、追加元は通常の第二対応を同時 $Z$ 倍に移し、
指定した第二対応を保たない。反転なし衝突では $T(c,1/c)$ の二乗が第二対応に自明となるには
$c=\pm1$ が必要で、全反転と合わせてやはり同じ四元群である。
両向きのせん断の作用が単射なので、全安定群の恒等成分も指定トーラスだけである。
軸交換で逆方向の流入は順方向と同一視できる。
滑らかな閉軌道横断面で順方向の接指標は重み1が4個、重み2が3個である。
定理\ref{thm:split-surface-full-fibers}の正規表示の一次核とせん断余空間は定階数の束となり、
四法線式の余核で測る二次障害は一般点で恒等零なので全域で恒等零である。
これにより全パラメータで同じ重み付き射影商と全次数正則性を得る。
横断面でのトーラス軌道境界は原点のみである。
追加単純表示では接次元61、アフィン軌道次元60で、一つの横断方向はパラメータを動かす
重み零の方向である（同じ軌道次元の曲線族が全て多重安定）。
従って正の横断座標はなく、全流入が閉軌道だけで尽きる。

二移動の持上げは旧式の $\kappa=-(u^2-1)(v^2-1)/(4uv)$ の平方根を使う。
余空間と核が束になったので、標数零の各補正対合の固有次元は被覆上で局所一定であり、
既出の一般値がそのまま全点の値になる。平方根の符号は重み1側の正負を交換するだけである。
非零正方向を固定する比較は指定トーラスを逆にできないので、可能な追加の基底比較は上記四元群に尽きる。
四成分の相違の証明も対合の交換則と固有次元によるため全域に適用できる。

残る一移動は $t=1+\sqrt2$、$h=1$ に取る。重み1,2の座標を $b,c$ とする。
補正 $X$ を法線数 $n$ でさらに $\iota^{n-2}$ 倍して対合に取り直す。
補正一移動 $S$ は四語を左巡回した後に全文字2,3だけのビットを反転し、
旧語の各位置・文字に順に次表の係数を掛け、全体を $t^{n-2}$ 倍する。
\[
\begin{pmatrix}
 \iota&1&\iota t^2&-t^2\\
 -\iota&-1&\iota&1\\
 -\iota/t^2&-1/t^2&-\iota&1\\
 \iota&1&\iota&1
\end{pmatrix}.
\]
これが背景を保つこと、$S^4=1$ であることを直接検算できる。横断座標の切断は
法線トーラスと $X,Z,S$ を含む群に同変に取れる。この群は分解トーラスの有限拡大で簡約的である。
従って作用は $b,c$ 上の線型作用で決まる。$g=1,Z,X,XZ$ とし、
$b\ne0$ なら $gSb=\lambda b,\ gSc=\lambda^2c$、$b=0$ なら $c$ が $gS$ の固有ベクトル
という条件が一移動の必要十分条件である。
\path{computations/verify_split_collision_shift.py} は両切断の接核、せん断商、
各作用を $\mathbb Q(\sqrt2,\iota)$ で取り、全ての $gS$ が $b$ 上に四乗根各一つを持つことを確認する。
$c$ 上の固有値は $g=1,Z$ で $1,1,-1$、$g=X,XZ$ で $1,\iota,-\iota$ である。
各固有部分空間の組の包含を四元群の全像で比較すると極大なものの代表は
$(g,\lambda)=(1,1),(1,\iota),(X,1),(X,\iota)$ と、
$b=0$, $g=X$ で $c$ が $\iota$ 固有の一点である。
その $b,c$ の許容次元の対は $(1,2),(1,1),(1,1),(1,0),(0,1)$ となり、主張の五成分を得る。
\end{proof}
